% from ega4-2-fr.tex
\oldpage[IV-2]{195}

\begin{env}[7.3.10]
\label{IV.2.s7.3.10-fr}
Nous dirons qu'une propriété $\boldsymbol P(Z,k)$ est du \emph{premier
type} si la propriété $\boldsymbol Q(A,k)$ qui sert à définir
$\boldsymbol P$ est une propriété $\boldsymbol R(A)$ \emph{ne faisant pas
intervenir $k$} et est vraie lorsque $A$ est un \emph{corps}; nous dirons
que $\boldsymbol P(Z,k)$ est du \emph{second type} si la propriété
$\boldsymbol Q(A,k)$ est de la forme
\begin{quote}
« Pour toute extension $k'$ de type fini de $k$ et tout point $z'$ de
$\operatorname{Spec}(A\otimes_k k')$ au-dessus du point fermé de
$\operatorname{Spec}(A)$, on a $\boldsymbol R(\mathscr O_{z'})$ »
\end{quote}
$\boldsymbol R(A)$ étant encore supposée vraie lorsque $A$ est un corps.
Il est clair que dans les exemples de (7.3.8), les propriétés (i) à (vii)
sont du premier type, les propriétés (iv$'$) à (vii$'$) du second type.

Cela étant, si l'on reprend les raisonnements de (6.6.1) et (6.8.3), on
constate que lorsque $\boldsymbol P$ est du premier ou du second type pour
une propriété $\boldsymbol R(A)$, \emph{les conditions
$(\mathrm P_{\mathrm I})$ et $(\mathrm P_{\mathrm{II}})$ sont
conséquences des conditions suivantes sur $\boldsymbol R$} :

$(\mathrm R_{\mathrm I})$ Soient $X$, $Y$ deux préschémas localement
noethériens, $f:X\to Y$ un morphisme \emph{régulier}; alors, pour tout
$x\in X$, la propriété $\boldsymbol R(\mathscr O_{f(x)})$ implique la
propriété $\boldsymbol R(\mathscr O_x)$.

$(\mathrm R_{\mathrm{II}})$ Soit $\varphi:A\to B$ un homomorphisme local
d'anneaux locaux noethériens faisant de $B$ un $A$-module \emph{plat};
alors $\boldsymbol R(B)$ implique $\boldsymbol R(A)$.

En outre, la propriété $(\mathrm P_{\mathrm{III}})$ résulte de ce que
$\boldsymbol R(k)$ est vraie par hypothèse pour tout corps $k$; enfin,
lorsque $\boldsymbol P$ est du second type, la condition
$(\mathrm P_{\mathrm{IV}})$ est conséquence de la définition de
$\boldsymbol P$ et de la transitivité des extensions de type fini.
\end{env}

\begin{remark}[7.3.11]
\label{IV.2.s7.3.11-fr}
Nous laissons au lecteur le soin de formuler la propriété $\boldsymbol R$
correspondant à chacun des exemples de (7.3.8), et de vérifier les
conditions $(\mathrm R_{\mathrm I})$ et $(\mathrm R_{\mathrm{II}})$ dans
chacun des cas, à l'aide des résultats du \S\,6. En fait, sauf pour
l'exemple (i) de (7.3.8), la propriété $\boldsymbol R$ vérifie dans les
autres cas la condition suivante :
\begin{enumerate}
\item[$(\mathrm R'_{\mathrm I})$] Soient $X$, $Y$ deux préschémas
localement noethériens, $f:X\to Y$ un $\boldsymbol P$-morphisme (pour la
propriété $\boldsymbol P$ du premier ou du second type définie par
$\boldsymbol R$); alors, pour tout $x\in X$, la propriété
$\boldsymbol R(\mathscr O_{f(x)})$ implique la propriété
$\boldsymbol R(\mathscr O_x)$.
\end{enumerate}
Les raisonnements de (6.6.1) et (6.8.3) prouvent que si
$\boldsymbol R$ vérifie les conditions $(\mathrm R'_{\mathrm I})$ et
$(\mathrm R_{\mathrm{II}})$, alors $\boldsymbol P$ vérifie les conditions
$(\mathrm P'_{\mathrm I})$ (7.3.4, (iii)) et
$(\mathrm P_{\mathrm{II}})$.
\end{remark}

\begin{proposition}[7.3.12]
\label{IV.2.s7.3.12-fr}
Soit $\boldsymbol P$ une propriété du premier ou du second type définie à
partir d'une propriété $\boldsymbol R$ vérifiant les conditions
$(\mathrm R_{\mathrm I})$ et $(\mathrm R_{\mathrm{II}})$ (resp.
$(\mathrm R'_{\mathrm I})$ et $(\mathrm R_{\mathrm{II}})$). Si pour tout
préschéma localement noethérien $X$, $U_{\boldsymbol R}(X)$ désigne
l'ensemble des $x\in X$ tels que $\boldsymbol R(\mathscr O_x)$ soit vraie,
alors, pour tout morphisme régulier (resp. tout
$\boldsymbol P$-morphisme) $f:X\to Y$ de préschémas localement
noethériens, on a
\begin{equation}
U_{\boldsymbol R}(X)=f^{-1}(U_{\boldsymbol R}(Y)).
\tag{7.3.12.1}
\label{IV.2.eq.7.3.12.1-fr}
\end{equation}
\end{proposition}

\begin{proof}
C'est une conséquence immédiate des définitions.
\end{proof}

\begin{env}[7.3.13]
\label{IV.2.s7.3.13-fr}
Étant donné un anneau semi-local noethérien $A$, nous appellerons
\emph{fibres formelles de $A$} les fibres du morphisme canonique
$f:\operatorname{Spec}(\widehat A)\to\operatorname{Spec}(A)$; pour tout
idéal premier $\mathfrak p$ de $A$, la fibre formelle en $\mathfrak p$ est
donc le schéma
$\operatorname{Spec}(\widehat A\otimes_A k(\mathfrak p))$; comme le
complété de l'anneau local $A/\mathfrak p$ est
$\widehat A/\mathfrak p\widehat A=(A/\mathfrak p)\otimes_A\widehat A$, on
voit que la fibre formelle de $A$ en $\mathfrak p$ est aussi \emph{la
fibre formelle de $A/\mathfrak p$ au point générique (0) de
$\operatorname{Spec}(A/\mathfrak p)$}.
\end{env}

